\chapter{第二十一章：Agentic Engineering 实战 ------ 大型移动端项目的 AI 工程化体系}

\begin{quote}
基于某大型移动端项目的真实实践，2026-06\\
\end{quote}

\bigskip\hrule\bigskip

\section{目录}

\begin{enumerate}[leftmargin=2em]
\item 为什么要做 Agentic Engineering
\item 整体架构一览
\item Agent 开发详解
\item Skill 开发详解
\item Hook 质量门禁
\item Rules 条件加载规则
\item Eval 评测体系
\item 多平台指令镜像
\item 设计模式与最佳实践
\item 从 0 到 1 创建你的第一个 Skill
\item 从 0 到 1 创建你的第一个 Agent
\item 反模式与踩坑记录
\item 附录：该项目现有 Agent/Skill 速查表
\end{enumerate}

\bigskip\hrule\bigskip

\section{1. 为什么要做 Agentic Engineering}

\subsection{背景}

该项目是某大厂移动端 AI 产品的核心 SDK，包含 120+ 个 Kotlin/Compose 模块和对应的 iOS SwiftUI 模块，嵌入多个宿主 App。团队日常面临：


\begin{itemize}[leftmargin=2em]
\item \textbf{重复性高}：加 feature flag、加 snapshot test、加 telemetry event 都有固定套路
\item \textbf{PR review 成本大}：跨 iOS/Android 双平台，review 规则多且容易遗漏
\item \textbf{值班 DRI 工作琐碎}：事件管理平台分诊、发版、周报都是机械流程
\item \textbf{质量事故有规律}：线上 regression 往往因为漏跑某个 review pass
\end{itemize}

\subsection{核心理念}

\begin{quote}
\textbf{Treat AI agents as team members with guardrails.}\\
给 agent 足够自主权去完成工程任务，但通过 hook、eval、review-pass 形成多层质量闭环。\\
\end{quote}

这不是"用 AI 写代码"那么简单------而是构建一个 \textbf{AI 工程师团队}，每个 agent 有明确职责、工具链、知识库和约束条件。


\bigskip\hrule\bigskip

\section{2. 整体架构一览}

\begin{codeblock}
.claude/                          <- Claude Code 的配置根目录
 agents/                       <- Agent 定义（角色 + 工具 + 知识库）
    Nova.agent.md             <- 自主编码 agent
    Arbiter.agent.md          <- PR 质量评审 agent
    DRI.agent.md              <- 值班管理 agent
    Arbiter/                  <- Arbiter 的知识库子目录
       evaluation-flow.md
       rubric-ios.md
       rubric-android.md
       scorecard-template.md
    DRI/
        config.json
        workflows/
            1-dashboard-monitoring.md
            2-incident-triage.md
            ...

 skills/                       <- 可组合能力模块
    pr-risk-evaluate/
       SKILL.md              <- 核心定义文件
       guardrails.md         <- 辅助知识
       evals/                <- 评测用例
    android-add-feature-flag/
       SKILL.md
    ...（40+ skills）

 hooks/                        <- 工具调用拦截器
    pr-review-gate.sh         <- 创建 PR 前检查 review pass
    pre-push-ci-gate.sh       <- 推送前检查 CI marker
    swift-lint-on-edit.sh     <- 编辑 Swift 后自动 lint
    mark-review.sh            <- 记录 review 状态

 rules/                        <- 条件加载的编码规范
    ios-development.md        <- 仅 ios/**/*.swift 时加载
    android-testing.md        <- 仅 android/**/test/** 时加载
    pr-format.md              <- 始终加载

 evals/                        <- Agent 评测脚本
    run-nova-eval.py
    verify-nova-eval.py

 resources/                    <- 共享知识文件
    kusto_telemetry_table_reference.md

 commands/                     <- 快捷命令
 settings.json                 <- 权限 + hook 绑定
\end{codeblock}

\subsection{数据流}

\begin{codeblock}
用户请求


     意图路由
  Agent 层    ->   Skill 层
  (角色决策)                      (能力执行)


        调用工具（MCP/CLI/编辑文件）


  Hook 层     <- 拦截/增强 ->   Rules 层
  (质量门禁)                      (编码规范)




  Eval 层      <- 事后验证 agent/skill 输出质量
  (评测闭环)

\end{codeblock}

\bigskip\hrule\bigskip

\section{3. Agent 开发详解}

\subsection{3.1 Agent 文件结构}

一个 agent 由 \textbf{一个 @@P0@@ 文件} + \textbf{可选的知识库子目录} 组成。


\begin{codeblock}
.claude/agents/
 MyAgent.agent.md          <- 主定义文件
 MyAgent/                  <- 知识库目录（可选）
     workflow-a.md
     workflow-b.md
     config.json
\end{codeblock}

\subsection{3.2 \texttt{.agent.md} 文件格式}

\begin{codeblock}
---
name: MyAgent
description: "一句话描述 agent 的职责和触发场景。"
argument-hint: "告诉用户可以传什么参数，如 'a PR number' or 'LOCAL_DIFF'"
tools:
  [
    "read",
    "edit",
    "search",
    "Bash(git:*)",
    "mcp__pm__*",          # MCP 工具用完全限定名
    "todos",
  ]
---

# MyAgent - 角色全称

## Instructions

@file:.claude/agents/MyAgent/workflow-a.md    <- 引用外部文件

## Overview

描述 agent 的核心职责...

## Trigger Phrases

| 触发词 |
| "do X"; "handle Y"; "run Z" |

## Step-by-Step Workflow

### Step 1: ...
### Step 2: ...
\end{codeblock}

\subsection{3.3 关键设计原则}

\subsubsection{原则 1：工具集最小化}

\begin{codeblock}
# [OK] 好：只声明需要的工具
tools: ["read", "search", "mcp__pm__repo_pull_request"]

# [X] 差：给所有权限
tools: ["*"]
\end{codeblock}

Agent 的 \texttt{tools} 字段定义了它能使用的工具白名单。遵循最小权限原则------DRI agent 不需要编辑文件的权限，Arbiter 不需要推代码的权限。


\subsubsection{原则 2：知识按需加载（Progressive Disclosure）}

\begin{codeblock}
## Intent Triage

| Intent | Workflow File to Read |
|--------|------------------------|
| 事件分诊 | `.claude/agents/DRI/workflows/2-incident-triage.md` |
| 发版    | `.claude/agents/DRI/workflows/3-release-cut.md` |
\end{codeblock}

\textbf{不要}在 agent 主文件里塞所有知识。DRI agent 有 5 个 workflow，但只在匹配意图后才加载对应的 workflow 文件。这样：

\begin{itemize}[leftmargin=2em]
\item 减少 context window 占用
\item 加载更快
\item 每个 workflow 可以独立迭代
\end{itemize}

\subsubsection{原则 3：结构化输出格式}

\begin{codeblock}
## Output Format

| Dimension | Score | Assessment |
|-----------|-------|------------|
| Scope     | X/25  | ...        |
| Semantic  | X/25  | ...        |
\end{codeblock}

给 agent 定义严格的输出模板。这样：

\begin{itemize}[leftmargin=2em]
\item 人类可预期输出结构
\item 下游 skill/hook 可以解析
\item eval 可以验证格式正确性
\end{itemize}

\subsubsection{原则 4：Todo-Driven 执行}

\begin{codeblock}
## CREATE TODO LIST FIRST

**BEFORE starting any coding work, you MUST create a todo list.**

### Required Todo Structure
1. Understand the requirement
2. Explore codebase
3. Plan implementation
4. Create feature branch
5. Implement changes
6. Add tests
7. Run quality checks
8. Create PR
\end{codeblock}

Nova agent 强制要求先创建 todo list 再开始编码。这提供了：

\begin{itemize}[leftmargin=2em]
\item 对用户的可见性（知道 agent 在做什么）
\item 进度追踪
\item 可中断/可恢复
\end{itemize}

\subsection{3.4 现有 Agent 案例分析}

\subsubsection{Nova（自主编码 agent）}

\begin{itemize}[leftmargin=2em]
\item \textbf{角色}：给定项目管理平台 ticket，端到端实现功能
\item \textbf{知识库}：40KB+ 的 coding-workflow.md（业界最详细的 agent workflow 之一）
\item \textbf{关键设计}：
\item 强制 todo list -> 每一步都有可见性
\item 调用 skill（如 \texttt{android-add-feature-flag}）而不是自己写脚手架代码
\item 完成后自动触发 \texttt{pr-risk-evaluate} 和 \texttt{review-pr}
\end{itemize}

\subsubsection{Arbiter（PR 质量评审 agent）}

\begin{itemize}[leftmargin=2em]
\item \textbf{角色}：独立评审 PR 质量，6 维度 5 级打分
\item \textbf{知识库}：平台特定 rubric（rubric-ios.md / rubric-android.md）
\item \textbf{关键设计}：
\item 和 Nova 形成 \textbf{"maker-checker"} 模式------Nova 写代码，Arbiter 审代码
\item 有 scorecard 模板确保输出一致
\item 支持 \texttt{--build-free} 模式（跳过构建，只做静态分析）
\end{itemize}

\subsubsection{DRI（值班管理 agent）}

\begin{itemize}[leftmargin=2em]
\item \textbf{角色}：IM 事件分诊 + 发版 + 周报
\item \textbf{MCP 工具链}：IM（事件管理）+ 日志查询平台 + 即时通讯（发消息）
\item \textbf{关键设计}：
\item 支持 \texttt{/loop 10m @DRI check dashboard} ------ 每 10 分钟自动巡检
\item 有 \texttt{config.json} 控制自动分诊行为
\item 工作流严格分离到 5 个独立文件
\end{itemize}

\bigskip\hrule\bigskip

\section{4. Skill 开发详解}

\subsection{4.1 Skill vs Agent 的区别}

| | Agent | Skill |

|--|-------|-------|

| \textbf{比喻} | 一个有角色的工程师 | 工程师掌握的一项技能 |

| \textbf{触发} | 用户直接对话 | 被 agent 或用户调用 |

| \textbf{状态} | 有持续的角色上下文 | 无状态，执行完即结束 |

| \textbf{复杂度} | 高（多步决策、意图路由） | 中（步骤明确的标准流程） |

| \textbf{文件} | \texttt{.agent.md} + 知识库目录 | \texttt{SKILL.md} + 辅助文件 + evals |


\subsection{4.2 Skill 文件结构}

\begin{codeblock}
.claude/skills/
 my-skill/
     SKILL.md              <- 核心定义（必须）
     helper-data.md        <- 辅助知识文件（可选）
     reference-rules.md    <- 参考规则（可选）
     evals/                <- 评测用例（强烈推荐，G6 guardrail 要求）
         eval-001.md
         eval-002.md
         ...
\end{codeblock}

\subsection{4.3 \texttt{SKILL.md} 文件格式}

\begin{codeblock}
---
name: my-skill
description: Does X for Y in the project. Use when asked to "do X",
             "handle Y", or "run Z for module W".
---

# My Skill

简短描述这个 skill 做什么。

## Prerequisites

- 列出前置条件
- 需要哪些工具/环境

## Step-by-Step Instructions

### 1. 第一步
具体指令...

### 2. 第二步
具体指令...

## Output Format

定义输出格式...
\end{codeblock}

\subsection{4.4 SKILL.md 写作规范}

根据该项目的 \texttt{skill-review} skill 中总结的 best practices：


\begingroup
\scriptsize
\setlength{\tabcolsep}{3pt}
\begin{longtable}{|p{0.31\textwidth}|p{0.31\textwidth}|p{0.31\textwidth}|}
\hline
\# & 规则 & 示例 \\ \hline
1 & \textbf{Description 用第三人称} & [OK] \texttt{"Adds feature flag to Android module"} [X] \texttt{"I can help add..."} \\ \hline
2 & \textbf{Description 包含触发词} & \texttt{"Use when asked to 'add feature flag' or 'create toggle'"} \\ \hline
3 & \textbf{Description <= 1024 字符} & 不含 XML tag \\ \hline
4 & \textbf{Name 格式} & 小写 + 连字符，<=64 字符，如 \texttt{android-add-feature-flag} \\ \hline
5 & \textbf{SKILL.md <= 500 行} & 超长内容拆到辅助文件 \\ \hline
6 & \textbf{引用文件只能一层深} & SKILL.md -> helper.md [OK]，SKILL.md -> a.md -> b.md [X] \\ \hline
7 & \textbf{不硬编码时间/版本} & [X] \texttt{"As of June 2026..."} \\ \hline
8 & \textbf{MCP 工具用完全限定名} & [OK] \texttt{mcp\_\_pm\_\_wit\_get\_work\_item} [X] \texttt{pm/get\_work\_item} \\ \hline
9 & \textbf{不假设工具存在} & 需要 pip 包就写 \texttt{pip install X} \\ \hline
10 & \textbf{用 Unix 路径} & 即使在 Windows 上也用 \texttt{/} \\ \hline
\end{longtable}
\endgroup

\subsection{4.5 Skill 分类}

该项目的 40+ skills 大致分为以下几类：


\subsubsection{脚手架类（Scaffolding）}
自动化常见的代码生成任务：

\begin{itemize}[leftmargin=2em]
\item \texttt{android-add-feature-flag} --- 加 feature flag
\item \texttt{android-add-snapshot-test} --- 加 Paparazzi snapshot test
\item \texttt{android-add-telemetry-event} --- 加遥测事件
\item \texttt{android-add-use-case} --- 加 UseCase
\item \texttt{android-add-new-module} --- 创建新 Gradle 模块
\end{itemize}

\subsubsection{质量检查类（Quality）}
确保代码/PR 质量：

\begin{itemize}[leftmargin=2em]
\item \texttt{pr-risk-evaluate} --- PR 风险评分（0-100，4 维度 + 14 条 guardrail）
\item \texttt{review-pr} --- PR 代码审查（规则 + swarm 并行）
\item \texttt{pr-quality-check} --- detekt + coverage 检查
\item \texttt{code-coverage} --- 覆盖率分析
\item \texttt{skill-review} --- 审查其他 skill 的质量
\end{itemize}

\subsubsection{流程自动化类（Workflow）}
自动化日常流程：

\begin{itemize}[leftmargin=2em]
\item \texttt{create-pr} --- PR 创建流程
\item \texttt{address-pr-comments} --- 处理 PR review 评论
\item \texttt{shiproom-update} --- 发版会议更新
\item \texttt{triage-business-chat-bugs} --- Bug 分诊
\end{itemize}

\subsubsection{调试类（Debug）}
辅助问题排查：

\begin{itemize}[leftmargin=2em]
\item \texttt{pm-pipeline-analyzer} --- Pipeline 失败分析
\item \texttt{cud-kusto} --- Kusto 查询辅助
\end{itemize}

\bigskip\hrule\bigskip

\section{5. Hook 质量门禁}

\subsection{5.1 Hook 是什么}

Hook 是在 AI agent 执行工具调用 \textbf{前后} 自动触发的 shell 脚本。它们是 agentic guardrails 的核心实现机制。


\subsection{5.2 配置方式}

在 \texttt{.claude/settings.json} 中定义：


\begin{codeblock}
{
  "hooks": {
    "PreToolUse": [
      {
        "matcher": "mcp__pm__repo_create_pull_request",
        "hooks": [
          {
            "type": "command",
            "command": "bash \".claude/hooks/pr-review-gate.sh\"",
            "timeout": 10,
            "statusMessage": "PR-review gate: checking..."
          }
        ]
      }
    ],
    "PostToolUse": [
      {
        "matcher": "Edit|Write",
        "hooks": [
          {
            "type": "command",
            "command": "bash \".claude/hooks/swift-lint-on-edit.sh\"",
            "timeout": 10
          }
        ]
      }
    ]
  }
}
\end{codeblock}

\subsection{5.3 Hook 分类}

\begingroup
\scriptsize
\setlength{\tabcolsep}{3pt}
\begin{longtable}{|p{0.19\textwidth}|p{0.19\textwidth}|p{0.19\textwidth}|p{0.19\textwidth}|p{0.19\textwidth}|}
\hline
时机 & Hook & 匹配 & 作用 &  \\ \hline
\textbf{创建 PR 前} & \texttt{pr-review-gate.sh} & \texttt{repo\_create\_pull\_request} & 检查是否跑过 review-pr / pr-risk-evaluate &  \\ \hline
\textbf{推送代码前} & \texttt{pre-push-ci-gate.sh} & \texttt{Bash(git push)} & 检查本地 CI 是否通过 &  \\ \hline
\textbf{编辑文件后} & \texttt{swift-lint-on-edit.sh} & `Edit\ & Write` & 自动对编辑的 Swift 文件跑 lint \\ \hline
\textbf{编辑文件后} & \texttt{claude-md-quality-check.sh} & `Edit\ & Write` & 检查是否修改了 CLAUDE.md \\ \hline
\textbf{编辑文件后} & \texttt{cowork-swift-strict-on-edit.sh} & `Edit\ & Write` & Cowork 项目严格检查 \\ \hline
\end{longtable}
\endgroup

\subsection{5.4 PR Review Gate 深度解析}

这是最重要的 hook，让我们详细看看 \texttt{pr-review-gate.sh} 的逻辑：


\begin{codeblock}
# 1. 判断是否在创建 PR
case "$TOOL" in
  *repo_create_pull_request*) is_pr_create=true ;;
esac

# 2. 计算 branch diff，找到改了哪些文件
DIFF=$(git diff "$RANGE" --name-only)

# 3. 根据改动类型决定需要哪些 review pass
need pr-risk        && missing+=("pr-risk")          # 每个 PR 都需要
echo "$DIFF" | grep -qE 'ios|android' &&             # 代码 PR 需要
  need review-pr    && missing+=("review-pr")
echo "$DIFF" | grep -qiE '\.(swift|kt|png)$' &&      # UI PR 需要
  need visual       && missing+=("visual")
echo "$DIFF" | grep -qE '\.claude/skills/' &&         # Skill PR 需要
  need skill-review && missing+=("skill-review")

# 4. 如果有缺失的 pass，发出提醒
echo "{\"additionalContext\":\"请先运行 $LIST\"}"
\end{codeblock}

关键机制：

\begin{itemize}[leftmargin=2em]
\item \texttt{mark-review.sh} 在每个 review pass 运行后写入 commit SHA 到 \texttt{.claude/.review-state/}
\item \texttt{pr-review-gate.sh} 检查当前 HEAD 是否有对应的标记
\item 如果有新 commit，标记过期，需要重新运行
\end{itemize}

\subsection{5.5 Hook 的 Advisory vs. Blocking 模式}

目前该项目的所有 hook 都是 \textbf{advisory}（建议模式）------输出 \texttt{additionalContext} 提醒 agent，但不阻止执行。


\begin{codeblock}
# Advisory（当前使用）
echo '{"additionalContext":"请先运行 review-pr"}'
exit 0

# Blocking（可选升级）
echo '{"hookSpecificOutput":{"permissionDecision":"deny","permissionDecisionReason":"..."}}'
\end{codeblock}

设计文档中明确标注了升级路径，但出于稳定性考虑暂时保持 advisory。


\bigskip\hrule\bigskip

\section{6. Rules 条件加载规则}

\subsection{6.1 规则定义}

Rules 是 \textbf{按文件路径条件加载} 的编码规范。只有当 agent 编辑匹配路径的文件时才会激活。


\begin{codeblock}
<!-- .claude/rules/ios-development.md -->
<!-- paths: ios/**/*.swift -->

# iOS Development Rules

- 使用 Swift Concurrency (async/await)，不用 Combine
- 使用 @Observable，不用 ObservableObject
- Composable 函数前缀用 PascalCase
- ...
\end{codeblock}

\subsection{6.2 该项目的 Rules 体系}

\begingroup
\scriptsize
\setlength{\tabcolsep}{3pt}
\begin{longtable}{|p{0.31\textwidth}|p{0.31\textwidth}|p{0.31\textwidth}|}
\hline
Rule 文件 & 加载条件 & 内容 \\ \hline
\texttt{ios-development.md} & 编辑 \texttt{ios/**/*.swift} & iOS 架构、CLI 命令、Swift 规范 \\ \hline
\texttt{ios-testing.md} & 编辑 \texttt{ios/**/Tests/**} & Snapshot test、单元测试模式 \\ \hline
\texttt{android-development.md} & 编辑 \texttt{android/**/*.kt} & Android 架构、Gradle、Kotlin 规范 \\ \hline
\texttt{android-testing.md} & 编辑 \texttt{android/**/test/**} & 单元测试、Compose test 模式 \\ \hline
\texttt{pr-format.md} & 始终加载 & PR 标题/描述格式 \\ \hline
\end{longtable}
\endgroup

\subsection{6.3 为什么不把所有规则放在一个文件里？}

\begin{itemize}[leftmargin=2em]
\item \textbf{Context window 效率}：iOS 开发者不需要加载 Android 规则
\item \textbf{维护性}：各平台团队独立维护自己的规则
\item \textbf{精确性}：测试规则只在编辑测试文件时加载，避免干扰正常编码
\end{itemize}

\bigskip\hrule\bigskip

\section{7. Eval 评测体系}

\subsection{7.1 为什么需要 Eval}

Agent 和 Skill 本质上是 \textbf{prompt engineering 的产物}。和代码一样，它们需要测试来保证质量。Eval 就是 agent/skill 的单元测试。


\subsection{7.2 Eval 文件格式}

\begin{codeblock}
<!-- evals/eval-001-simple-pr.md -->

## Input
Evaluate PR #XXXXX

## Expected Behavior
- Should fetch PR metadata from PM平台
- Should compute a risk score 0-100
- Should identify platform as iOS
- Score should be in LOW RISK range (0-30)
- Output should include Dimension Breakdown table
\end{codeblock}

\subsection{7.3 G6 Guardrail --- Eval 的强制执行}

该项目的 \texttt{pr-risk-evaluate} skill 中定义了 G6 guardrail：


\begin{quote}
\textbf{G6（agentic system）}: 当 PR 修改了 \texttt{.claude/skills/}、\texttt{.claude/agents/}、\texttt{.claude/commands/} 下的文件时触发。\\
- 检查被修改 skill 目录下的 eval 文件数量\\
- 要求 \textbf{>= 50 个 eval} 且 \textbf{100\% pass rate}\\
- 不满足则 floor score 设为 70（HIGH RISK）\\
\end{quote}

这形成了一个自我进化的闭环：

\begin{codeblock}
修改 skill -> G6 要求 eval -> 必须先写/更新 eval -> eval 不过不让创建 PR
\end{codeblock}

\subsection{7.4 Eval 运行机制}

\begin{codeblock}
# .claude/evals/run-nova-eval.py
# 自动运行 Nova agent 的评测用例，验证输出质量
\end{codeblock}

\bigskip\hrule\bigskip

\section{8. 多平台指令镜像}

\subsection{8.1 问题}

该项目同时使用两个 AI 编码工具：

\begin{itemize}[leftmargin=2em]
\item \textbf{Claude Code} --- 读取 \texttt{.claude/rules/*.md}
\item \textbf{GitHub Copilot CLI} --- 读取 \texttt{.github/instructions/*.instructions.md}
\end{itemize}

两个工具都不会自动读取对方的目录。


\subsection{8.2 解决方案：双向镜像}

\begin{codeblock}
.claude/rules/ios-development.md
     手动保持同步
.github/instructions/ios-development.instructions.md
\end{codeblock}

每个文件头部都有注释指向镜像文件：


\begin{codeblock}
<!-- Mirror of .claude/rules/ios-development.md. When you change one, update the other. -->
\end{codeblock}

在 \texttt{AGENTS.md} 中明确规定了同步规则：


\begin{quote}
When you modify a file in either folder, update its counterpart in the other folder.\\
Filenames pair 1:1. To add a new rule, create both files in the same change.\\
\end{quote}

\subsection{8.3 为什么不用一个共享文件？}

\begin{itemize}[leftmargin=2em]
\item Claude Code 的 \texttt{paths:} 用 YAML 列表格式
\item Copilot CLI 的 \texttt{applyTo:} 用逗号分隔字符串
\item 两个工具都不支持 markdown 链接跨目录引用
\item 接受重复，换取两个工具都能正常工作
\end{itemize}

\bigskip\hrule\bigskip

\section{9. 设计模式与最佳实践}

\subsection{模式 1：Maker-Checker（制造者-检查者）}

\begin{codeblock}
Nova (maker) -> 写代码 -> Arbiter (checker) -> 评审质量
\end{codeblock}

两个 agent 形成分权制衡。Nova 写的 PR 必须经过 Arbiter 评审，就像人类工程师写代码要过 code review。


\subsection{模式 2：Skill 组合（Composition）}

\begin{codeblock}
Nova agent
   调用 android-add-feature-flag skill
   调用 android-add-snapshot-test skill
   自己写业务逻辑
   调用 pr-quality-check skill
   调用 create-pr skill
\end{codeblock}

Agent 不重复发明轮子------复用 skill 完成标准化子任务。


\subsection{模式 3：Guard-then-Act（先守后行）}

\begin{codeblock}
创建 PR (action)
    ^
pr-review-gate.sh (guard)
    检查 review-pr ?
    检查 pr-risk-evaluate ?
    检查 design-visual-audit ?
\end{codeblock}

Hook 在 action 前拦截，确保前置条件满足。


\subsection{模式 4：从 Regression 学习（Incident-Driven Guardrails）}

该项目的每条 guardrail 都关联真实的线上事故：


\begin{codeblock}
G6  (agentic system)   <- 防止未经测试的 skill 变更
G12 (iPad coverage)    <- 来自 commit xxxxxxx 的 iPad regression
G13 (flag collapse)    <- 来自 #XXXXX 的 feature flag bug
G14 (custom toolbar)   <- 来自 #XXXXX/#XXXXX 的 iOS 26 glass bug
\end{codeblock}

这不是凭空设计的规则------每一条都是从血的教训中提炼的。


\subsection{模式 5：Intent-Driven Workflow Loading（意图驱动加载）}

\begin{codeblock}
## Intent Triage

| Intent         | Workflow to Load                |
|----------------|---------------------------------|
| 事件分诊       | workflows/2-incident-triage.md  |
| 发版           | workflows/3-release-cut.md      |
| 周报           | workflows/4-weekly-report.md    |
\end{codeblock}

不预加载所有 workflow------匹配意图后才加载，节省 context window。


\subsection{模式 6：Swarm Review（群体审查）}

\begin{codeblock}
review-pr skill
   规则审查（review-rules-*.md）
   Swarm 并行审查
      Breaker（找 bug）
      Exploiter（找安全问题）
      Inspector（找架构问题）
   合并去重 -> 发 PM平台 评论
\end{codeblock}

单一 reviewer 的盲区通过多角色并行审查来弥补。


\bigskip\hrule\bigskip

\section{10. 从 0 到 1 创建你的第一个 Skill}

\subsection{示例：创建一个 \texttt{android-add-string-resource} skill}

\subsubsection{Step 1：创建目录结构}

\begin{codeblock}
mkdir -p .claude/skills/android-add-string-resource/evals
\end{codeblock}

\subsubsection{Step 2：编写 SKILL.md}

\begin{codeblock}
---
name: android-add-string-resource
description: Adds a new string resource to the Android project with
             proper localization setup. Use when asked to "add a string",
             "create string resource", or "add localized text".
---

# Add Android String Resource

Adds a new string entry to the project string resource system, including
the base `strings.xml` and the `StringResource` Kotlin constant.

## Prerequisites

- Working in the `project/android/` directory
- Module name where the string is needed

## Step-by-Step Instructions

### 1. Identify the Target Module
Determine which module needs the string (e.g., `ui`, `chat`, `business-chat`).

### 2. Add to strings.xml
Add the string entry to `<module>/src/main/res/values/strings.xml`:
\end{codeblock}
<string name="app\_<feature>\_<description>">English text</string>

\begin{codeblock}

### 3. Add StringResource Constant
Add a constant to the module's `StringResource` provider:
\end{codeblock}
val FeatureDescription = StringResource("app\_<feature>\_<description>")

\begin{codeblock}

### 4. Verify
Run: `project-cli gradle :<module>:assembleDebug`
\end{codeblock}

\subsubsection{Step 3：添加评测用例}

\begin{codeblock}
<!-- evals/eval-001-basic-add.md -->
## Input
Add a string "Cancel" for the chat input module.

## Expected Behavior
- Creates entry in chat module's strings.xml
- Adds StringResource constant
- Uses naming convention app_chat_cancel
\end{codeblock}

\subsubsection{Step 4：测试你的 skill}

\begin{codeblock}
# 在 Claude Code 中
/android-add-string-resource Add a "Send" button label to the ui module
\end{codeblock}

\bigskip\hrule\bigskip

\section{11. 从 0 到 1 创建你的第一个 Agent}

\subsection{示例：创建一个 \texttt{Reviewer} agent}

\subsubsection{Step 1：创建文件}

\begin{codeblock}
touch .claude/agents/Reviewer.agent.md
mkdir -p .claude/agents/Reviewer
\end{codeblock}

\subsubsection{Step 2：编写 agent 定义}

\begin{codeblock}
---
name: Reviewer
description: "Reviews Android Compose code for performance anti-patterns.
              Use when asked to 'check compose performance' or
              'review recomposition'."
tools:
  [
    "read",
    "search",
    "mcp__pm__repo_pull_request_thread_write",
  ]
---

# Reviewer - Compose Performance Agent

## Overview
Reviews Jetpack Compose code for common performance issues:
recomposition, unstable types, missing keys, etc.

## Trigger Phrases
| "check compose performance"; "review recomposition"; "find unstable composables" |

## Workflow

### Step 1: Identify Compose Files
Search for `@Composable` in changed files.

### Step 2: Check Anti-Patterns
For each file, check:
- Missing `@Stable` / `@Immutable` on data classes passed to composables
- Lambda allocations inside composition
- Missing `key()` in `LazyColumn` items
- Unnecessary recomposition triggers

### Step 3: Report Findings
Post findings as PM平台 PR comments with severity and fix suggestions.

## Rules Reference
@file:.claude/agents/Reviewer/compose-rules.md
\end{codeblock}

\subsubsection{Step 3：添加知识库}

\begin{codeblock}
<!-- .claude/agents/Reviewer/compose-rules.md -->
# Compose Performance Rules

## Rule 1: Unstable Parameters
Data classes passed to @Composable must be @Immutable or @Stable.
...

## Rule 2: Lambda Allocation
Avoid creating lambdas inside composition scope.
...
\end{codeblock}

\bigskip\hrule\bigskip

\section{12. 反模式与踩坑记录}

\subsection{[X] 反模式 1：Agent 文件过大}

\begin{codeblock}
<!-- 40KB 的 agent 文件 = 每次对话都消耗大量 context -->
\end{codeblock}

\textbf{解决}：拆分到子目录，按需加载。Nova 的 40KB workflow 是一个已知问题。


\subsection{[X] 反模式 2：Skill 之间链式引用}

\begin{codeblock}
SKILL.md -> a.md -> b.md -> c.md    <- Claude 不会跟超过一层
\end{codeblock}

\textbf{解决}：所有引用文件只能从 SKILL.md 直接引用（一层深度）。


\subsection{[X] 反模式 3：硬编码版本号}

\begin{codeblock}
[X] "Use Compose 1.5.4"
[OK] "Use the version defined in gradle/libs.versions.toml"
\end{codeblock}

\subsection{[X] 反模式 4：忘记双平台同步}

修改了 \texttt{.claude/rules/ios-development.md} 但忘了更新 \texttt{.github/instructions/ios-development.instructions.md}。


\textbf{解决}：在 AGENTS.md 中明确规定，每个文件头部有镜像提示。


\subsection{[X] 反模式 5：Guardrail 没关联事故}

\begin{codeblock}
[X] G15: Check for magic numbers
[OK] G13: New flag collapses a conditional branch --- grounded in #XXXXX
\end{codeblock}

没有事故支撑的 guardrail 缺乏说服力，容易被忽略。


\subsection{[X] 反模式 6：重命名 Paparazzi 测试方法但不删旧快照}

\begin{codeblock}
# CLAUDE.md 中明确警告：
# DO NOT rename Paparazzi test methods without deleting the old snapshot image
\end{codeblock}

\bigskip\hrule\bigskip

\section{13. 附录：该项目现有 Agent/Skill 速查表}

\subsection{Agents}

\begingroup
\scriptsize
\setlength{\tabcolsep}{3pt}
\begin{longtable}{|p{0.23\textwidth}|p{0.23\textwidth}|p{0.23\textwidth}|p{0.23\textwidth}|}
\hline
Agent & 角色 & 知识库 & MCP 集成 \\ \hline
\textbf{Nova} & 自主编码 & coding-workflow.md (40KB) & PM平台 \\ \hline
\textbf{Arbiter} & PR 质量评审 & rubric-ios/android.md, scorecard & PM平台 \\ \hline
\textbf{DRI} & 值班管理 & 5 个 workflow, config.json & IM, Kusto, 即时通讯 \\ \hline
\textbf{Atlas} & 架构规划 & iPad checklist, extension points & --- \\ \hline
\textbf{Jing} & --- & --- & --- \\ \hline
\textbf{Scorpion} & --- & --- & --- \\ \hline
\textbf{Touchstone} & --- & --- & --- \\ \hline
\textbf{Debuggy} & iOS 调试 & --- & Appium, Jupyter \\ \hline
\textbf{Hao} & --- & orchestration-workflow.md & --- \\ \hline
\end{longtable}
\endgroup

\subsection{Skills（按类别）}

\begingroup
\scriptsize
\setlength{\tabcolsep}{3pt}
\begin{longtable}{|p{0.47\textwidth}|p{0.47\textwidth}|}
\hline
类别 & Skills \\ \hline
\textbf{脚手架} & android-add-feature-flag, android-add-snapshot-test, android-add-telemetry-event, android-add-use-case, android-add-new-module, android-add-provider, android-add-chat-setting, android-add-debug-setting, android-add-feature-toggle \\ \hline
\textbf{质量} & pr-risk-evaluate, review-pr, pr-quality-check, code-coverage, skill-review, review-dev-spec, review-partner-pr, project-pr-ultrareview \\ \hline
\textbf{流程} & create-pr, address-pr-comments, shiproom-update, triage-business-chat-bugs, pages-strings \\ \hline
\textbf{构建/运行} & android-build-and-run, android-build-and-run-on-mac, android-build-and-run-on-windows, android-start-emulator-on-windows, setup \\ \hline
\textbf{测试} & automated-ui-test, android-appium-setup, copilot-notebooks-android-appium, appium-setup, android-pre-push-patterns \\ \hline
\textbf{调试/分析} & pm-pipeline-analyzer, cud-kusto, exp-scorecard-analyzer \\ \hline
\textbf{审查} & accessibility, design-visual-audit (implied) \\ \hline
\textbf{Cowork} & cowork-architecture-review, cowork-flaky-test-stabilizer, cowork-pr-review, cowork-snapshot-update, cowork-swift-concurrency, android-cowork-project \\ \hline
\textbf{验证} & pacman-request-validation \\ \hline
\end{longtable}
\endgroup

\subsection{Hooks}

\begingroup
\scriptsize
\setlength{\tabcolsep}{3pt}
\begin{longtable}{|p{0.31\textwidth}|p{0.31\textwidth}|p{0.31\textwidth}|}
\hline
Hook & 时机 & 作用 \\ \hline
\texttt{pr-review-gate.sh} & 创建 PR 前 & 检查 review pass \\ \hline
\texttt{pre-push-ci-gate.sh} & git push 前 & 检查 CI marker \\ \hline
\texttt{swift-lint-on-edit.sh} & 编辑后 & Swift lint \\ \hline
\texttt{cowork-swift-strict-on-edit.sh} & 编辑后 & Cowork 严格 lint \\ \hline
\texttt{claude-md-quality-check.sh} & 编辑后 & CLAUDE.md 质量 \\ \hline
\texttt{mark-review.sh} & review 后 & 记录 review 状态 \\ \hline
\texttt{write-green-ci-marker.sh} & CI 通过后 & 写 CI 标记 \\ \hline
\end{longtable}
\endgroup

\bigskip\hrule\bigskip

\emph{本文档基于某大型移动端项目代码库 2026-06 版本整理。Agent/Skill 体系持续演进中。}
